\chapter*{Liste des sigles et des abréviations}
\def\coldeux{\relax}
\def\linestretch{3pt}

\begin{twocolumnlist}{.2\textwidth}{.7\textwidth}
 BvK & Born-von-Karman, se référant aux conditions de bord de Born-von-Karman ou le cristal de Born-von-Karman\\
 BZ & Zone de Brillouin, de l'anglais \textit{Brillouin Zone}\\
 DFT & Théorie de la fonctionnelle de la densité, de l'anglais
      \textit{Density Functional Theory}\\
 DFPT & Théorie de la fonctionnelle de la densité perturbative, de l'anglais \textit{Density Functional Perturbation Theory}\\
 FFT & Transformée de Fourier rapide, de l'anglais \textit{Fast Fourier Transform}\\
 GGA & Approximation du gradient généralisé, de l'anglais \textit{Generalized Gradient Approximation}\\
 DMFT & Théorie du champ moyen dynamique, de l'anglais \textit{Dynamical Mean Field Theory}\\
 HK & Hohenberg et Kohn, se référant aux théorèmes de Hohenberg et Kohn\\
 IBZ & Zone de Brillouin irréductible, de l'anglais \textit{Irreducible Brillouin Zone}\\
 KB & Kleinman et Bylander, se référant à la forme séparable de Kleinman et Bylander pour les pseudo-potentiels à conservation de norme\\
 KS & Kohn-Sham, se référant à la méthode de Kohn-Sham\\
 LDA & Approximation de la densité locale, de l'anglais \textit{Local Density Approximation}\\
 MLWF & Fonction de Wannier maximalement localisée, de l'anglais \textit{Maximally Localized Wannier Fonction}\\
 MP & Monkhorst-Pack, se référant à la grille de points-$\kv$ de Monkhorst-Pack\\
 NC & Conservation de la norme, de l'anglais \textit{Norm-Conserving}\\
 NSCF & Méthode du champ non-auto-cohérent, de l'anglais \textit{Non Self Consistent Field}\\
 ONCV & Conservation de norme optimisé de Vanderbilt, se référant aux pseudo-potentiels à conservation de norme optimisé de Vanderbilt (de l'anglais \textit{Optimized Norm-Conserving Vanderbilt})\\
 PBE & Perdew-Burke-Ernzerhof, se référant à la fonctionnelle d'échange-corrélation Perdew-Burke-Ernzerhof\\
 SCF & Méthode du champ auto-cohérent, de l'anglais \textit{Self Consistent Field}\\
 WF & Fonction de Wannier\\
 XC & Échange-corrélation, se référant à la fonctionnelle ou l'énergie d'échange-corrélation\\
\end{twocolumnlist}